\documentclass[letterpaper,11pt]{article}
\usepackage[top=0.8in, bottom=0.8in, left=0.9in, right=0.9in]{geometry}
\usepackage[hidelinks]{hyperref}
\usepackage{lmodern}
\usepackage{parskip}
\usepackage{microtype}
\setlength{\parskip}{6pt}
\setlength{\parindent}{0pt}
\begin{document}
\begin{flushleft}
\textbf{\large Ajay Venkatesh} \\
Brooklyn, NY \\
+1 516 585 9013 \\
\href{mailto:av3855@nyu.edu}{av3855@nyu.edu}
\end{flushleft}
\vspace{10pt}
May 21, 2026
\vspace{6pt}

Hiring Manager \\
Odoo

\vspace{10pt}

Dear Hiring Manager,

My name is Ajay Venkatesh. I'm finishing my M.S. in Computer Engineering at NYU, with production software experience at PwC in Bengaluru, a summer SDE internship at Amazon, and ongoing on-campus work at NYU's Novel AI Technologies. I'm writing about the Full Stack Engineer role on the Upgrade Services team.

The role's framing of diving into custom Python and JS code written years ago by different teams, then managing complex database migrations with minimal disruption, is exactly the kind of pragmatic engineering I find most interesting. At PwC I spent the past few years modernizing legacy onboarding systems, owning a distributed backend on \textbf{Microsoft Azure} integrating with \textbf{Microsoft CRM} over \textbf{SQL Server}, with parallelized batch ingestion that materially cut execution time. Understanding what historical code is actually doing before refactoring it is a skill that consulting-style delivery teaches well.

On full-stack work: at NYU's Novel AI Technologies I shipped a customer-facing analytics platform end-to-end with \textbf{React, Node.js, REST APIs}, deployed via \textbf{Docker} on \textbf{EC2}, with Stripe billing and row-level security for paying clients. At Amazon I built a \textbf{React 18} frontend with optimized API integration alongside backend services in \textbf{Java} with low-level design patterns, all shipped under code-review discipline.

On database migration and complex data work: at NYU I engineered \textbf{ETL pipelines} transforming \textbf{NoSQL} into \textbf{PostgreSQL} with concurrency control and locking mechanisms, processing hundreds of thousands of records. At PwC I optimized data ingestion with parallel processing on multi-thousand-row Excel feeds, cutting execution time from hours to minutes. The discipline of validating outputs against source systems carries over directly to upgrade migrations.

Stack-wise: \textbf{Python}, \textbf{JavaScript / TypeScript}, \textbf{SQL}, \textbf{PostgreSQL}, plus \textbf{Object-oriented programming}, \textbf{GitHub} workflows, and ERP-adjacent enterprise systems from PwC. Comfortable across legacy-system rehabilitation, refactoring, and migration testing.

What pulls me toward Odoo is the engineering discipline. Open-source ERP for thousands of companies runs on careful refactoring and migration work where mistakes translate into real client data loss, and the chance to learn from custom code written by different teams over years is exactly the kind of compounding engineering education I'm looking for.

I'd welcome the chance to discuss further. Thanks for considering my application.

\vspace{12pt}
Sincerely, \\
Ajay Venkatesh

\end{document}